\documentclass{article}

\usepackage{amsmath,amssymb}
\usepackage{xurl}
\usepackage[hidelinks]{hyperref}

% Shared commands for notes in docs/paper-gaps/.

\DeclareUnicodeCharacter{2080}{\ifmmode{}_{0}\else\textsubscript{0}\fi}
\DeclareUnicodeCharacter{2081}{\ifmmode{}_{1}\else\textsubscript{1}\fi}
\DeclareUnicodeCharacter{2082}{\ifmmode{}_{2}\else\textsubscript{2}\fi}
\DeclareUnicodeCharacter{2099}{\ifmmode{}_{n}\else\textsubscript{n}\fi}

\newcommand{\Lean}{\textsc{Lean}}
\ExplSyntaxOn
\NewDocumentCommand{\leanid}{m}{
  \ifmmode
    \textsf{\detokenize{#1}}
  \else
    \group_begin:
    \urlstyle{sf}
    \tl_set:Nn \l_tmpa_tl {#1}
    \tl_replace_all:Nnn \l_tmpa_tl {\_} {_}
    \exp_args:NV \path \l_tmpa_tl
    \group_end:
  \fi
}
\ExplSyntaxOff
\providecommand{\tr}{\operatorname{tr}}
\newcommand{\tnleanIssueBase}{https://github.com/LionSR/TNLean/issues/}
\newcommand{\tnleanPrBase}{https://github.com/LionSR/TNLean/pull/}
\newcommand{\tnleanDocBase}{https://sirui-lu.com/TNLean/docs/}
\newcommand{\ghissue}[1]{\href{\tnleanIssueBase#1}{GitHub issue \##1}}
\newcommand{\ghpr}[1]{\href{\tnleanPrBase#1}{PR \##1}}
\newcommand{\doclink}[1]{\href{\tnleanDocBase#1.html}{\path{#1}}}

% Dirac notation and the identity operator, as in blueprint/src/macros/common.tex.
% \providecommand keeps note-local definitions authoritative.
\usepackage{dsfont}
\providecommand{\ket}[1]{|#1\rangle}
\providecommand{\bra}[1]{\langle #1|}
\providecommand{\braket}[2]{\langle #1 | #2 \rangle}
\providecommand{\ketbra}[2]{|#1\rangle\!\langle #2|}
\providecommand{\Id}{\mathds{1}}
\providecommand{\one}{\mathds{1}}
\providecommand{\C}{\mathbb{C}}
\providecommand{\MN}[1]{M_{#1}(\C)}
\providecommand{\MD}{M_D(\C)}

% External referencing for paper sources.  The build helper writes sanitized
% auxiliary files under build/paper-aux/ and excludes duplicated source labels.
\usepackage{xr}

% Import a generated paper auxiliary file with a caller-supplied label prefix.
% The candidate paths support builds from docs/paper-gaps/, the repository root,
% and build/paper-aux/ itself.
\newcommand{\importpaperaux}[2]{%
  \IfFileExists{../../build/paper-aux/#2.aux}{%
    \externaldocument[#1]{../../build/paper-aux/#2}%
  }{%
    \IfFileExists{build/paper-aux/#2.aux}{%
      \externaldocument[#1]{build/paper-aux/#2}%
    }{%
      \IfFileExists{#2.aux}{%
        \externaldocument[#1]{#2}%
      }{}%
    }%
  }%
}

% General external reference macros
\newcommand{\paperref}[2]{\ref{#1-#2}}
\newcommand{\papereqref}[2]{(\ref{#1-#2})}

% CPSV16 source (1606.00608 / MPDO-22-12-17-2.tex) external document and shortcuts
\importpaperaux{cpsv16-}{1606.00608/MPDO-22-12-17-2-tnlean}
\newcommand{\cpsvref}[1]{\paperref{cpsv16}{#1}}
\newcommand{\cpsveqref}[1]{\papereqref{cpsv16}{#1}}

% Verdict marker: \gapnote{<kind>}{<status>}. Kind is the policy
% classification (clarification, local-correction, scope-restriction,
% unfaithful, false-source, open-gap); status is open, wip (elimination
% underway), resolved, or historical. Machine-read by the site index;
% prints a verdict line.
\newcommand{\gapnote}[2]{\par\noindent\textbf{Verdict:} #1 (#2).\par}


\title{The L-Symbol Relation Imposed by Time Reversal}
\author{}
\date{2026-09-26}

\begin{document}
\maketitle
\gapnote{local-correction}{resolved}

\section{The printed statement}

Ref.~\cite{GarreRubio2023MPOSym}, subsection ``Interplay between MPO and
time reversal symmetry'' (lines 865--958 of the source), considers a matrix
product state with blocks $x$ that is symmetric under a group matrix product
operator symmetry $(O_g)_{g\in G}$ and under an antiunitary time reversal
$\mathcal T=v^{\otimes n}\mathcal K$ commuting with every $O_g$. Write
$V_{g,x}$ for the action tensor mapping block $x$ to block $g\cdot x$ and
$L^x_{g,h}$ for the L-symbols defined by
\begin{equation}
  V_{g,h\cdot x}V_{h,x}=L^x_{g,h}\,W_{g,h}V_{gh,x},
  \label{eq:f1}
\end{equation}
the equation \texttt{F1group} of the source, written schematically with the
fusion tensor $W_{g,h}$. Comparing the local actions of $\mathcal TO_g$ and
$O_g\mathcal T$, the source finds nonzero scalars $\gamma_{g,x}$ with
$S_{g\cdot x}^{-1}V_{g,x}S_x=\gamma_{g,x}V_{g,x}^*$, and from them the
display at lines 935--952,
\begin{equation}
  \gamma_{h,x}\gamma_{g,h\cdot x}\,(V_{g,h\cdot x}V_{h,x})^*
  =L^x_{g,h}\,\gamma_{gh,x}\,(W_{g,h}V_{gh,x})^*.
  \label{eq:display}
\end{equation}
Equation \texttt{TRSF1group} at line 954 then states
\begin{equation}
  |L^x_{g,h}|^2=\frac{\gamma_{gh,x}}{\gamma_{h,x}\gamma_{g,h\cdot x}},
  \label{eq:printed}
\end{equation}
and the text concludes that the L-symbols squared are trivial because they
can be written as a gauge transformation.

\section{The printed relation fails}

Rescaling the action tensors $V_{g,x}\mapsto c_{g,x}V_{g,x}$ multiplies
$L^x_{g,h}$ by $c_{g,h\cdot x}c_{h,x}/c_{gh,x}$ and $\gamma_{g,x}$ by the
unimodular number $c_{g,x}/\overline{c_{g,x}}$. The left side of
(\ref{eq:printed}) changes by a positive factor that can take any value,
while the right side changes by a unimodular factor, so (\ref{eq:printed})
cannot hold in every gauge.

An explicit instance with an on-site symmetry, where $W_{g,h}$ is trivial
and there is one block, is the following. Let $G=\mathbb Z_2\times\mathbb
Z_2=\{e,a,b,ab\}$ and let $A$ be the real matrix product state tensor with
components $\mathbb 1,X,iY,Z$, where $X,Y,Z$ are the Pauli matrices. It is
invariant under $\mathcal T=\mathcal K$ with $S=\mathbb 1$, and under the
real diagonal on-site symmetry $u_g$ whose virtual action is conjugation by
the action tensors
\[
  V_e=\mathbb 1,\qquad V_a=X,\qquad V_b=Z,\qquad V_{ab}=Y .
\]
Since $V_e,V_a,V_b$ are real and $Y^*=-Y$, the scalars $\gamma$ are
$\gamma_e=\gamma_a=\gamma_b=1$ and $\gamma_{ab}=-1$. From $XZ=-iY$ one has
$L_{a,b}=-i$. Then $|L_{a,b}|^2=1$, while
$\gamma_{ab}/(\gamma_b\gamma_a)=-1$. The numerical check with
$V_b=e^{i\theta}Z$ gives $|L_{a,b}|^2=1$ against the ratio
$-e^{-2i\theta}$.

\section{The corrected statement}

Taking the complex conjugate of (\ref{eq:f1}) gives
$(V_{g,h\cdot x}V_{h,x})^*=\overline{L^x_{g,h}}\,(W_{g,h}V_{gh,x})^*$.
Substituting into (\ref{eq:display}) and cancelling the nonzero tensor
$(W_{g,h}V_{gh,x})^*$ yields
\begin{equation}
  \frac{\overline{L^x_{g,h}}}{L^x_{g,h}}
  =\frac{\gamma_{gh,x}}{\gamma_{h,x}\gamma_{g,h\cdot x}} .
  \label{eq:corrected}
\end{equation}
In the example, $\overline{L_{a,b}}/L_{a,b}=i/(-i)=-1$, in agreement with
(\ref{eq:corrected}). Since $\overline L/L=(|L|/L)^2$, the relation
(\ref{eq:corrected}) says that the square of the phase $L/|L|$ is a gauge
transformation of the trivial L-symbol; this is the correct reading of the
statement that the L-symbols squared are trivial. The modulus $|L|$ is not
constrained. The printed (\ref{eq:printed}) replaces the ratio
$\overline L/L$ by $\overline LL$.

For a general matrix product operator symmetry the fusion tensors also pick
up time-reversal phases $\beta_{g,h}$, as at lines 873--889 of the source,
and the gauge-covariant form of (\ref{eq:corrected}) is that
$\overline L$ equals the joint gauge transform of $L$ by $(\beta,\gamma)$,
jointly with $\overline\omega$ equal to the gauge transform of $\omega$ by
$\beta$. Under these two relations a single joint gauge makes both $\omega$
and $L$ real, by taking unimodular square roots of the phases of $\beta$ and
$\gamma$.

\section{Formalized statement}

The formalization proves, at the level of scalars, the cancellation step
from (\ref{eq:display}) and the conjugate of (\ref{eq:f1}) to
(\ref{eq:corrected}); the statement that the squared phase of $L$ is the
gauge ratio; the reality of the anomaly under lines 873--891; and the joint
reality of $\omega$ and $L$.\footnote{Formal declarations:
\leanid{TNLean.Algebra.LSymbol.star\_apply\_eq\_gauge\_one\_apply},
\leanid{TNLean.Algebra.LSymbol.phase\_sq\_eq\_of\_gauge\_eq\_star},
\leanid{TNLean.Algebra.ScalarThreeCochain.exists\_cohomologousTo\_star\_eq\_self}
and \leanid{TNLean.Algebra.LSymbol.exists\_gauge\_star\_eq\_self} in
\doclink{TNLean/Algebra/ScalarThreeCocycleTimeReversal}.}
The antiunitary symmetry itself, and the derivation of $\gamma$ and $\beta$
from $\mathcal T$ commuting with $O_g$, are not formalized.

\textbf{Verdict: resolved local correction.} The printed
\texttt{TRSF1group} is false; the source's own preceding display gives
(\ref{eq:corrected}), which is the relation used in the formalization, and
the source's conclusion that the L-symbols squared are trivial holds in
that corrected form.

\bibliographystyle{alpha}
\bibliography{references}

\end{document}
